\section{Conclusion}\label{sec:conclusion}

We set out to measure how much of what providers could report they choose not to. The coded record returns the number and the reason to distrust it: $+11.4$ points, significant against zero, never more than 1.3 points from its label-free counterpart. Dating the instruments bounds the endowment from above, the step the disclosure literature has long wanted, but the same dates decide where a model sits inside a benchmark's coverage; the variable identifying what could have been held also moves the measure of what was withheld. Two things follow. A measure of standing used to price silence must be defended as free of the timing that decides disclosability; reweighting is a partial defence at best. And growing the drop record, 16 realised drops against a ceiling of 37 pairs, is a reporting-standard question with two levers: a scale estimated jointly across benchmarks \citep{epoch2026rosetta}, which here shrinks the placebo gap without closing it, and carrying a benchmark set across a family's releases, which removes the empty-endowment excuse and narrows the \citet{dye1985} pooling without new disclosure.
